% Tabela po vrsticah: najprej samo glava, nato X, Y in Z.
\begin{frame}{Odkrivanje tabele po vrsticah}
    Včasih pride prav, da tabelo odkrivamo postopoma po vrsticah.
    \begin{center}
        \begin{tabular}{c|cccc}
            Oznaka & A & B & C & D \\ \hline
            \only<2->{X & 1 & 2 & 3 & 4 \\}
            \only<3->{Y & 3 & 4 & 5 & 6 \\}
            \only<4->{Z & 5 & 6 & 7 & 8 \\}
        \end{tabular}
    \end{center}
\end{frame}

% Tabela po stolpcih: vsaka kolona ima svoj predpis za prikaz.
\begin{frame}{Odkrivanje tabele po stolpcih}
    Tabelo lahko odkrivamo tudi po stolpcih, čeprav ni najlažje.
    \begin{center}
        \begin{tabular}{c|>{\onslide<2->}c<{\onslide} >{\onslide<3->}c<{\onslide} >{\onslide<4->}c<{\onslide} >{\onslide<5->}c<{\onslide}}
            Oznaka & A & B & C & D \\ \hline
            X & 1 & 2 & 3 & 4 \\
            Y & 3 & 4 & 5 & 6 \\
            Z & 5 & 6 & 7 & 8
        \end{tabular}
    \end{center}
\end{frame}
